\originalchapter{22.02 作业三:隧穿与质量差}
题源: Spring 2012 Problem Set 3. 源 PDF 发布日为3月4日, 截止年份印为2018, 保留该元数据错误, 不改变 Spring 2012 课程归属. 本章解答为编者独立补写.
\begin{exercise}\label{ex:2202-3-1}
原题 1. 质子能量 $E=15$ MeV, 入射高度 $V_H=30$ MeV、宽 $L=10$ fm 方势垒. (a) 用 $T\simeq4\e^{-2\kappa L}$ 估计; (b) 用式 \eqref{eq:barrier} 精确式比较; (c) 近似何时有效, 本题是否满足? 证明厚势垒且 $E=V_H/2$ 时退化为近似式. (d) $E\ne V_H/2$ 而 $E<V_H$ 时怎样变化?
\end{exercise}
\begin{solution}
取 $m_pc^2=938.272$ MeV, $\hbar c=197$ MeV fm, $\kappa\simeq0.852\,\mathrm{fm}^{-1}$, $\kappa L\simeq8.52$. 近似 $T=1.60\times10^{-7}$, 精确值 $1/[1+\sinh^2(\kappa L)]=\operatorname{sech}^2(\kappa L)$, 与近似的相对差约 $8\times10^{-8}$. 因 $\sinh^2x\simeq\e^{2x}/4$, 题中前因子是 4. 一般能量的前因子为 $16E(V_H-E)/V_H^2$, 同时 E 增加使 $\kappa$ 减小, 透射概率趋增; 靠近势垒顶需用完整式, 厚势垒近似不再一致.
\end{solution}
\begin{exercise}\label{ex:2202-3-2}
原题 2. 中子从左射向半无限台阶, $V=0$ ($x<0$)、$V_H=2E$ ($x>0$). (a) 原题要求找区域 II 的``中子流''降至入射流 $1/4$ 的 x. (b) 对 $m_nc^2=939$ MeV、E=10 MeV 求数值.
\end{exercise}
\begin{solution}
题设的禁阻区只有衰减波, 实际概率流为零, 所以不存在流为 $1/4$ 的位置. 若原意是波强度相对入射密度降至 $1/4$, 则 $k=\kappa$, 边界 $|t|^2=2$, 方程 $2\e^{-2\kappa x}=1/4$ 给 $x=\log8/(2\kappa)\simeq1.495$ fm. 若是相对界面处密度降至 $1/4$, 则改为 $x=\log4/(2\kappa)$. 必须说明参照和密度/流的区别, 不能给非零禁阻区流.
\end{solution}
\begin{exercise}\label{ex:2202-3-3}
原题 3. ${}^{224}$Ra 的 $\alpha$ 半衰期3.66天. ${}^{224}$Ra$\to{}^{210}$Pb+$^{14}$C 能同 $\alpha$ 模式竞争吗? (a) 只看 Q 值; (b) 估计比较半衰期; (c) 查实验分支比.
\end{exercise}
\begin{solution}
AME2020 给 $Q_\alpha=5.789$ MeV, $Q_{14C}=30.534$ MeV, 两个都放能, 单靠 Q 无法排除任一通道. 对每通道用 $R=1.25(A_d^{1/3}+A_c^{1/3})$ fm、$C=1.439965Z_dZ_c$ MeV fm、$\mu=M_dM_c/(M_d+M_c)$, 第 \ref{ch:tunnel} 章 WKB 指数可复算势垒穿透和撞击频率. 更大的 cluster 电荷和约化质量使势垒指数大; 还需 cluster 预形成概率, 因而仅用 Q 最大判断竞争性是不对的.

把两种预形成概率都人为取 1 时, WKB 指数分别为 $59.21,68.67$, 撞击频率分别为 $8.85\times10^{20},1.02\times10^{21}$ s$^{-1}$, 半衰期约 $0.47$ 天与 $14.4$ 年, 供比较模型敏感性; 它们不能当作实测寿命. NUBASE2020 的 $^{224}$Ra 条目给 $^{14}$C 分支 $4.0\times10^{-9}\%$, 即总衰变概率约 $4.0\times10^{-11}$. 因 $\alpha$ 几乎占全部, $^{14}$C 部分半衰期约 $3.66\,\mathrm d/(4\times10^{-11})\simeq2.5\times10^8$ 年. 实验上完全不能竞争. 核数据检索原题 NNDC 入口本次返回服务错误, 数值以可取得的官方 NUBASE2020 表核对, 不声称已读取失败网页.
\end{solution}
\begin{exercise}\label{ex:2202-3-4}
原题 4. 用原子质量求 $^{211}$Po 和 $^{210}$Po 的 $\alpha$ 衰变子核、Q、两产物动能和速度.
\end{exercise}
\begin{solution}
子核分别是 $^{207}$Pb、$^{206}$Pb. AME2020 Q 分别 $7.59456,5.40753$ MeV. 非相对论反冲分配 $T_\alpha=QM_d/(M_d+m_\alpha)$, $T_d=Qm_\alpha/(M_d+m_\alpha)$. 得前者 $T_\alpha\simeq7.4505,T_d\simeq0.1441$ MeV, 后者 $5.3045,0.1030$ MeV. 速度分别由 $v/c=\sqrt{2T/(mc^2)}$ 算, 前者约 $0.0632,0.00122$, 后者 $0.0533,0.00104$. 若核质量与原子质量的电子细节所要求精度更高, 需补电子结合修正.
\end{solution}
\begin{exercise}\label{ex:2202-3-5}
原题 5. $^{213}$At 衰变到 $^{209}$Bi, 半衰期125 ns, 发射 $\alpha$ 动能9080 keV. 由此求 $^{213}$At 原子质量.
\end{exercise}
\begin{solution}
$Q=T_\alpha(1+m_\alpha/M_d)\simeq9.254$ MeV. 用 $M({}^{209}\mathrm{Bi})=208.980398599$ u、$M({}^4\mathrm{He})=4.00260325413$ u, 得 $M({}^{213}\mathrm{At})\simeq212.9929363$ u. 原题元素英文印作 Attinium, 实际符号 At 是砹 (astatine). 半衰期不参与这一质量差计算, 不能由寿命直接替代能量守恒.
\end{solution}
